(1) 相空间中，电子态密度贡献 $p_eE_e$，中微子态密度与动量相乘给 $E_\nu p_\nu c$。当 $m_\nu=0$，$p_\nu c=E_\nu=Q-E_e$，所以 $d\Gamma/dE_e\propto p_eE_e(Q-E_e)^2$。若电子质量不能忽略，就不可把 $p_ec$ 直接替换成 $E_e$。

(2) 电子也超相对论时，谱形 $f(E)\propto E^2(Q-E)^2$。计算归一积分 $\int_0^QE^2(Q-E)^2dE=Q^5\int_0^1x^2(1-x)^2dx=Q^5/30$。因此 $\boxed{f(E)=30E^2(Q-E)^2/Q^5}$，其在 $E=0$ 与 $E=Q$ 消失，对 $E\leftrightarrow Q-E$ 对称。

(3) 由对称性 $\langle E\rangle=Q/2$。令 $x=E/Q$，是Beta$(3,3)$ 分布，$\langle x^2\rangle=3\cdot4/(6\cdot7)=2/7$，所以 $\operatorname{Var}(E)=Q^2(2/7-1/4)=Q^2/28$。未归一积分随 $Q^5$ 缩放，因此固定弱相互作用矩阵元且粒子都取超相对论时总率 $\Gamma\propto Q^5$；真实核衰变还有电子质量、库仑Fermi因子和核结构影响。

(4) 质量非零时 $p_\nu c=\sqrt{E_\nu^2-m_\nu^2c^4}$，所以 $d\Gamma/dE\propto E^2(Q-E)\sqrt{(Q-E)^2-m_\nu^2c^4}\,\Theta(Q-E-m_\nu c^2)$。电子谱截止于 $\boxed{E_{\max}=Q-m_\nu c^2}$。令 $\delta=E_{\max}-E\downarrow0$，根号 $\sim\sqrt{2m_\nu c^2\delta}$，表现为平方根端点；与无质量时 $(Q-E)^2$ 近端点幂律不同，是实验质量拟合的基础。
